% !TEX root = ../main.tex
%
% \section 生成大型项目这一章的一级标题，并让本文件中的 subsection 都归到该章编号之下。
\section{大型项目：拆分、命令与排错}

% learningbox 中的 \texttt{\textbackslash input} 用等宽字体显示命令名，\textbackslash 用来打印反斜杠本身。
\begin{learningbox}{学习目标 / Learning goals}
学完本章后，你将能把长文档拆成多个章节，判断 \texttt{\textbackslash input} 和相对路径的起点，
定义不会意外覆盖已有命令的新命令，辨认常见辅助文件，并按固定顺序定位编译问题。
\end{learningbox}

% 目录树放在 codeblock 中，竖线、反引号和文件名会按原样显示，不会被当成 LaTeX 命令解释。
\subsection{为什么拆分文件}

文档变长后，把所有内容放进一个 \texttt{main.tex} 会让搜索、审阅和多人协作越来越困难。
更稳妥的做法是让根文件集中保存文档类、宏包、全局样式、章节顺序和参考文献输出，
再让每个章节文件只负责一个清楚的主题。这样可以单独定位内容，又不会复制全局配置。

下面是本教程完成后的项目树。它位于惰性源码环境中，只展示目录关系，不会让 LaTeX
尝试载入尚未讲解的文件。

\begin{codeblock}
latex-tutorial/
|-- .vscode/
|   `-- settings.json
|-- figures/
|   |-- example-figure.pdf
|   `-- example-figure.tex
|-- sections/
|   |-- 01-getting-started.tex
|   |-- 02-document-structure.tex
|   |-- 03-text-formatting.tex
|   |-- 04-tables-and-figures.tex
|   |-- 05-mathematics.tex
|   |-- 06-theorems-and-references.tex
|   |-- 07-large-projects.tex
|   `-- 08-beamer.tex
|-- main.tex
|-- main.pdf
|-- beamer-demo.tex
|-- beamer-demo.pdf
|-- references.bib
`-- README.md
\end{codeblock}

% \input{路径} 把目标 .tex 内容插入当前位置；路径可省略扩展名，但章节文件仍需通过 main.tex 获得导言区配置。
\subsection{input 和主文件}

根文件按阅读顺序载入章节；\verb|\input| 会把目标文件的内容插入当前位置，通常可以省略
\texttt{.tex} 扩展名。例如，\texttt{main.tex} 中的前两章写作：

\begin{codeblock}
% !TEX root = ../main.tex
%
% \section 创建一级标题，LaTeX 会自动编号，并把标题文字写入目录。
\section{入门：从源文件到 PDF}

% learningbox 是主文件定义的自定义环境；紧跟在环境名后的花括号会作为彩色框标题。
\begin{learningbox}{学习目标 / Learning goals}
学完本章后，你将理解 TeX、LaTeX、XeLaTeX、Biber 和 LaTeXmk 这五种工具的关系，
并能在 VS Code 中编译一个最小中文文档。
\end{learningbox}

% \subsection 创建 section 下面的二级标题，编号和目录层级都会由 LaTeX 自动处理。
\subsection{为什么使用 LaTeX}

LaTeX 把“写什么”和“长什么样”分开：你在纯文本源文件中标记标题、段落、公式、
图片和引用，再由编译工具统一生成 PDF。这种工作方式特别适合结构较长、公式较多、
需要交叉引用或多人使用同一版式的文档。源文件可以逐行阅读、搜索和比较，版式规则则
集中写在主文件的导言区，因而修改一处就能保持全文一致。

% 这条 \subsection 命令开始工具说明小节；标题中的顿号只是普通文字，不影响命令语法。
\subsection{TeX、LaTeX、XeLaTeX、Biber 与 LaTeXmk}

\noindent\begin{tabularx}{\textwidth}{>{\bfseries}l X X}
\toprule
工具 & 作用 & 本项目中的位置 \\
\midrule
TeX & 最底层的排版系统，提供命令展开、断行和分页等基础能力。 & 其他工具共同依赖的排版基础。 \\
LaTeX & 建立在 TeX 之上的文档系统，提供文档类、章节、列表和交叉引用。 & \texttt{main.tex} 与各章节所使用的标记语言。 \\
XeLaTeX & 能直接处理 Unicode 和系统字体的 LaTeX 编译引擎。 & 负责把本项目的中文源文件编译成 PDF。 \\
Biber & 读取 \texttt{.bib} 数据库并生成参考文献所需信息。 & 处理根目录中的 \texttt{references.bib}。 \\
LaTeXmk & 根据依赖关系自动重复运行 XeLaTeX，并在需要时调用 Biber。 & LaTeX Workshop 构建配方调用的自动化工具。 \\
\bottomrule
\end{tabularx}

% 最小文档只需要 \documentclass 和一对 document 环境；其他宏包与设置都可按需要逐步添加。
\subsection{第一个最小文档}

把下面内容保存为一个 \texttt{.tex} 文件，并使用 XeLaTeX 编译：

\begin{codeblock}
\documentclass[UTF8]{ctexart}

\begin{document}
你好，LaTeX！
\end{document}
\end{codeblock}


% resultbox 是主文件预先定义的无参数环境；\raggedright 只在盒子内部取消两端对齐，避免窄框出现过大的字间距。
\begin{resultbox}
\raggedright
生成的 PDF 中只会出现“你好，LaTeX！”。\texttt{\textbackslash documentclass} 以及
\texttt{\textbackslash begin\{document\}} 之前的导言区用于配置文档，它们本身不会被打印到 PDF。
\end{resultbox}

% 本小节标题由 \subsection 生成；下面的操作步骤属于正文，会按当前章节样式连续排版。
\subsection{在 VS Code 中编译和预览}

% 下一段使用 \texttt 显示命令名和快捷键；该命令只改变花括号内文字的字体，不会执行其中的内容。
本项目已经启用自动保存（auto save）和保存后自动构建（build on save）。停止输入约 1 秒后，
VS Code 会自动保存当前源文件；这次保存会触发默认的第一条配方
\texttt{latexmk (XeLaTeX)}，由 LaTeXmk 调用 XeLaTeX 更新 PDF。若想立即手工构建，
仍可按 \texttt{Command+Option+B}。

\begin{enumerate}
    \item 打开项目文件夹（open the project folder）。
    \item 打开 \texttt{main.tex}（open \texttt{main.tex}）。
    \item 修改一小处，停止输入约 1 秒，等待自动保存和默认 XeLaTeX 配方完成构建。
    \item 按 \texttt{Command+Option+V}，在 VS Code 右侧打开 PDF 预览。
    \item 把光标放在一行源码中，按 \texttt{Command+Option+J}，从源码定位到 PDF 对应位置。
    \item 在 PDF 中双击一段文字，返回左侧对应的源码位置。
\end{enumerate}

% SyncTeX 由编译参数 synctex=1 生成位置映射；源码到 PDF 和 PDF 回源码都读取同一个 .synctex.gz 文件。
按 \texttt{Command+Option+J} 的“源码到 PDF”和双击 PDF 的“PDF 回到源码”合称
SyncTeX（源码与 PDF 同步定位）。前者也叫正向定位（forward search），后者也叫反向定位
（inverse search）；它们让你不必在长源码和长 PDF 中重复滚动查找同一段内容。

% warningbox 会把 begin 与 end 之间的内容放进橙色提示框；环境边界必须成对且名称完全相同。
\begin{warningbox}
\texttt{.tex} 文件是可以用编辑器打开和修改的纯文本；PDF 是编译后生成的输出文件。
要修改内容，应编辑 \texttt{.tex} 源文件并重新编译，而不是直接修改 PDF。
\end{warningbox}

% codeblock 按原样显示内部字符，因此其中故意缺少右花括号的 \textbf 只会被打印，不会作为真正命令执行。
\begin{errorbox}
下面的命令缺少一个右花括号：
\begin{codeblock}
\textbf{这行少了右花括号
\end{codeblock}
编译器随后可能报告很多连锁错误。请先找到并修复日志中第一条有用的错误信息，
然后重新编译；后面的错误往往会随之消失。
\end{errorbox}

% itemize 创建无序列表；可选参数 label={$\square$} 把每个 \item 前的默认圆点改为空方框。
\subsection{本章检查清单}

\begin{itemize}[label={$\square$}]
    \item 我能指出 \texttt{\textbackslash documentclass} 用来选择文档类型和选项。
    \item 我能写出从 \texttt{\textbackslash begin\{document\}} 到 \texttt{\textbackslash end\{document\}} 的文档环境。
    \item 我知道本项目的根文件是 \texttt{main.tex}。
    \item 我知道停止输入约 1 秒会自动保存，并触发默认的 XeLaTeX 配方。
    \item 我会使用 \texttt{Command+Option+B} 立即手工编译。
    \item 我会使用 \texttt{Command+Option+V} 在右侧打开 PDF 预览。
    \item 我会使用 \texttt{Command+Option+J} 从源码定位到 PDF。
    \item 我会双击 PDF 返回源码，并知道这两个方向合称 SyncTeX。
\end{itemize}

% !TEX root = ../main.tex
%
% \section 的花括号中是要显示的一级标题；编号、目录条目和标题样式都由文档类统一生成。
\section{文档结构：命令、环境与页面组成}

% learningbox 接收一个标题参数，并把环境内部的多段正文统一放进可分页的蓝色框中。
\begin{learningbox}{学习目标 / Learning goals}
学完本章后，你将能区分命令、参数和环境，理解文档类、导言区与宏包的职责，
并能组织标题、摘要、目录、章节层次和页面设置。
\end{learningbox}

% \subsection 表示二级标题；它从属于当前 section，并会自动获得类似“2.1”的编号。
\subsection{命令、参数与环境}

LaTeX 命令通常以反斜杠开始。方括号中的可选参数用来微调行为，花括号中的必选参数
则告诉命令要处理什么内容。下面是两种通用写法；它们位于源码框中，只作为示意，
不会被当前文档执行。

% codeblock 会原样打印示例：方括号通常放可省略的选项，花括号放必须提供的参数，环境则用同名 begin 和 end 包围内容。
\begin{codeblock}
\command[可选参数]{必选参数}

\begin{环境名}
环境中的内容
\end{环境名}
\end{codeblock}

在真实命令 \texttt{\textbackslash textbf\{粗体文字\}} 中，
\texttt{textbf} 是命令名，花括号中的文字是必选参数；实际执行后得到
\textbf{粗体文字}。花括号负责划定参数范围，本身不会显示在 PDF 中。

% quote 是实际执行的引用环境，会缩进其中的段落；它与源码框里的示意环境不同，必须由同名的 begin 和 end 正确闭合。
\begin{quote}
环境像一间有明确入口和出口的房间：\texttt{\textbackslash begin} 打开房间，
\texttt{\textbackslash end} 关闭房间，中间的内容按照该环境的规则排版。
\end{quote}

% errorbox 显示错误说明，内部再嵌套 codeblock；源码框中的错误列表不会运行，所以可以安全展示不匹配的环境名。
\begin{errorbox}
下面的列表用 \texttt{itemize} 开始，却错误地用 \texttt{enumerate} 结束：
\begin{codeblock}
\begin{itemize}
    \item 第一项
\end{enumerate}
\end{codeblock}
每个环境的开始名称和结束名称必须完全相同；这里应把
\texttt{\textbackslash end\{enumerate\}} 改成
\texttt{\textbackslash end\{itemize\}}，否则编译器会报告环境不匹配。
\end{errorbox}

% \documentclass 必须写在主文件导言区开头；方括号是类选项，最后一对花括号中的 ctexart 是文档类名称。
\subsection{文档类和文档选项}

\texttt{\textbackslash documentclass} 通常是主文件中第一条有效命令，而且一篇文档只选择
一个文档类。这里的 \texttt{ctexart} 是适合中文短文和教程的文章类；
\texttt{UTF8}、\texttt{12pt} 和 \texttt{a4paper} 分别指定编码、基础字号和纸张大小。

% 这个 codeblock 展示两类选项：documentclass 的选项控制整篇文档，usepackage 的选项只传给 geometry 宏包。
\begin{codeblock}
\documentclass[UTF8,12pt,a4paper]{ctexart}
\usepackage[left=2.6cm,right=2.6cm]{geometry}
\end{codeblock}

类选项写在文档类名之前的方括号中；若省略某个可选项，文档类会使用对应的默认值。
不同文档类支持的选项可能不同，因此不要把某个类的专用选项机械复制到另一个类中。

% 导言区位于 \documentclass 与 \begin{document} 之间，适合放宏包、全局设置和自定义命令，不放普通正文。
\subsection{导言区与宏包}

从 \texttt{\textbackslash documentclass} 之后到
\texttt{\textbackslash begin\{document\}} 之前的区域叫作导言区。
它适合加载宏包、定义颜色和命令、设置页边距，以及填写标题信息；导言区的配置通常不会
作为普通正文直接打印出来，却会影响整篇文档。

\texttt{\textbackslash usepackage} 用来加载宏包。例如，\texttt{geometry} 管理页面尺寸和
页边距，本教程还在主文件中加载了数学、图片、表格、列表、颜色和超链接等宏包。
把共享配置放在 \texttt{main.tex} 中，可以避免每个章节重复设置。

% warningbox 是主文件定义的提示环境；其中的 \texttt 把 LaTeX 命令以等宽字体显示，\textbackslash 专门打印反斜杠字符。
\begin{warningbox}
导言区用于“设置怎样排版”，\texttt{document} 环境用于“写要排版的内容”。
普通段落、列表和章节标题应放在 \texttt{\textbackslash begin\{document\}} 之后。
\end{warningbox}

% 标题信息先在导言区保存，再由正文中的 \maketitle 输出；abstract 和 tableofcontents 必须位于 document 环境内部。
\subsection{标题、摘要与目录}

标题、作者和日期通常在导言区声明；进入 \texttt{document} 环境后，
\texttt{\textbackslash maketitle} 才把这些信息排版出来。摘要由
\texttt{abstract} 环境包围，\texttt{\textbackslash tableofcontents} 根据带编号的章节命令
生成目录。

% 示例中的 \today 不接参数，会在每次编译时产生当天日期；\maketitle 必须放在标题、作者和日期定义之后。
\begin{codeblock}
\title{文档标题}
\author{作者姓名}
\date{\today}

\begin{document}
\maketitle
\begin{abstract}
这里是摘要。
\end{abstract}
\tableofcontents
\end{document}
\end{codeblock}

\texttt{\textbackslash today} 会在编译时生成当天日期。目录依赖辅助文件记录章节与页码，
所以刚加入或移动标题后，通常需要再次编译；本项目的 LaTeXmk 会自动完成所需轮次。

% 不带星号的 \subsection 会同时生成可见标题、自动编号和目录记录，通常不需要手写章节数字。
\subsection{章节层次与自动编号}

LaTeX 使用 \texttt{\textbackslash section}、\texttt{\textbackslash subsection} 和
\texttt{\textbackslash subsubsection} 表示从高到低的章节层次。本小节就是一个实际执行的
\texttt{subsection}：LaTeX 自动给它编号，并把它加入目录，而无需手工填写数字或页码。

% 命令名后的星号关闭自动编号和目录写入；花括号中的文字仍会作为小标题显示在 PDF 中。
\subsubsection*{无编号的小标题示例}

这行标题由实际的 \texttt{\textbackslash subsubsection*} 命令生成。带星号的标题
\textbf{没有编号，也不会自动加入目录}。如果确实需要让无编号标题出现在目录中，
应再使用专门的目录命令，而不是手写一个看似正确的页码。

% resultbox 用固定标题显示编译结果；其中修改章节结构后，LaTeXmk 会通过多轮编译自动更新编号、目录和页码。
\begin{resultbox}
使用章节命令后，新增、删除或移动小节时，LaTeX 会重新计算编号、目录和页码。
因此标题文字中不应手写“2.1”之类的序号。
\end{resultbox}

% 本小节中的 \clearpage 会先输出等待排版的图片和表格再换页；\pagestyle{plain} 选择只显示页码的基础页面样式。
\subsection{页面设置}

纸张大小可以由文档类选项指定，页边距则常交给 \texttt{geometry} 宏包统一管理。
本教程在主文件中使用 A4 纸，并分别设置上下左右页边距；这样各章节共享同一版心，
不需要在章节文件里重复调用宏包。

需要从新页开始时，可以使用 \texttt{\textbackslash clearpage}。它不仅换页，还会先安排
尚未排出的浮动图片和表格。\texttt{\textbackslash pagestyle\{plain\}} 则使用带页码的简洁页眉页脚样式。

% 这里再次使用 learningbox，但传入“本章检查”作为标题；同一个环境可在不同位置反复调用。
\begin{learningbox}{本章检查 / Chapter check}
我能解释命令参数与环境边界，指出导言区和正文的分界，按正确顺序组织标题、摘要与目录，
并说明普通标题与星号标题在编号和目录行为上的差别。
\end{learningbox}

\end{codeblock}

章节文件不是独立文档：它们不重复 \verb|\documentclass|、宏包或
\verb|document| 环境。即使当前正在编辑 \texttt{sections/07-large-projects.tex}，
也应把 \texttt{main.tex} 作为根文件编译。文件开头的编辑器提示
\verb|% !TEX root = ../main.tex| 能帮助 LaTeX Workshop 从子文件回到正确入口；
真正的构建结果仍由根文件中的章节顺序和全局配置决定。

% \newcommand 定义新命令；[1] 表示接收一个参数，定义体中的 #1 会替换为调用 \importantterm{...} 时提供的内容。
\subsection{自定义命令}

重复出现的语义和样式可以封装成命令。\verb|[1]| 表示命令接收一个参数，定义中的
\verb|#1| 会被调用时花括号里的内容替换。下面两行是活动内容，会实际参与本教程编译：

\newcommand{\importantterm}[1]{\textcolor{TutorialBlue}{\textbf{#1}}}
这里突出显示一个 \importantterm{重要术语}。

命令必须在第一次使用之前定义；全篇共用的命令通常放在根文件导言区。
LaTeX 还会检查新命令的名称是否已经存在。如果新名称与已有命令重名，
它会报告命令已定义，而不会静默覆盖原行为。只有确实要修改已知命令时，
才使用 \texttt{\textbackslash renewcommand}，并且应先确认原命令的参数形式。

% \path 会按路径样式显示斜杠等字符；这里的相对路径统一从实际编译入口 main.tex 所在目录开始计算。
\subsection{相对路径}

本项目从 \texttt{main.tex} 所在的 \texttt{latex-tutorial/} 目录构建。
根文件中的路径都从这个目录开始解释，例如：
\begin{itemize}[leftmargin=2.6em]
    \item \path{sections/07-large-projects} 指向章节源文件；
    \item \path{figures/example-figure.pdf} 指向图片；
    \item \path{references.bib} 指向文献数据库。
\end{itemize}
不要因为 \verb|\input| 命令写在某个子文件中，就自动把路径起点改成该子文件所在目录。

\begin{warningbox}
许多构建环境区分文件名和目录名的大小写，源码拼写必须与磁盘完全一致。
若从子文件启动了错误的编译入口，正确路径也会失效；此时先确认根文件，再检查路径。
\end{warningbox}

% description 的可选参数调整说明列表版式；每个 \item[扩展名] 把方括号内容作为术语标签，后面正文解释其用途。
\subsection{常见辅助文件}

编译会在源文件旁生成若干中间文件。它们记录跨轮次信息，通常不应手工编辑：

\begin{description}[style=nextline,leftmargin=3.2cm,font=\ttfamily]
    \item[.aux] 保存标签、引用和其他需要下一轮读取的交叉引用信息。
    \item[.log] 保存引擎的完整构建记录，是定位第一条有效错误的主要依据。
    \item[.toc] 保存目录条目，下一轮编译时由 \verb|\tableofcontents| 读取。
    \item[.out] 常由 \texttt{hyperref} 写入书签和链接相关信息。
    \item[.bcf] 把引用和排序等控制信息交给 Biber。
    \item[.bbl] 保存 Biber 根据 \texttt{.bib} 数据生成、供 LaTeX 读取的文献内容。
    \item[.blg] 保存 Biber 的处理日志和书目诊断信息。
    \item[.synctex.gz] 保存源码行与 PDF 位置的映射，支持正向和反向跳转。
    \item[.pdf] 是最终排版输出；修改内容后应由源码重新生成，而不是直接编辑它。
\end{description}

源文件 \texttt{.tex}、文献库 \texttt{.bib} 和图片资源是项目输入，不能当作辅助文件删除。
辅助文件虽然可以重建，但在正常增量构建中能帮助 LaTeXmk 判断依赖并减少不必要的工作。

% \small 临时缩小后面的表格；tabularx 的 X 列会自动分配剩余宽度，\arraybackslash 恢复单元格末尾使用双反斜杠的能力。
\subsection{固定排错顺序}

编译器可能在一次失败后输出许多连锁信息。先处理最早一条与自己源码有关的红色错误，
再重新构建，往往比从日志末尾倒查更有效。常见现象可以按下表缩小范围：

\small
\noindent\begin{tabularx}{\textwidth}{>{\ttfamily\raggedright\arraybackslash}p{0.25\textwidth} >{\raggedright\arraybackslash}X >{\raggedright\arraybackslash}X}
\toprule
\normalfont\bfseries 现象 & \bfseries 首先检查 & \bfseries 常见原因 \\
\midrule
Undefined control sequence & 命令拼写与宏包 & 命令拼错、命令尚未定义或缺少提供它的宏包。 \\
File not found & 路径与大小写 & 根文件不对、相对层级错误或资源没有放进项目。 \\
Missing \} inserted & 相关行之前的花括号 & 命令参数少了右花括号，错误位置可能晚于真正缺失处。 \\
Environment ended & \texttt{begin/end} 名称 & 环境名不一致、嵌套顺序错误或漏写结束命令。 \\
undefined references & 标签名与构建轮次 & 标签拼错，或修改编号后还没有完成足够的编译轮次。 \\
empty bibliography & 引用键、数据源与 Biber & 正文没有有效引用、键不匹配、文献库未载入或 Biber 未运行。 \\
\bottomrule
\end{tabularx}
\normalsize

% enumerate 自动生成顺序编号；nosep 由 enumitem 宏包提供，用来压缩各条目之间的额外竖直间距。
无论具体错误名称是什么，都按下面顺序逐项检查，并在每次只改一个原因后重新构建：

\begin{enumerate}[nosep]
    \item Read the first useful red error.（阅读第一条有用的红色错误。）
    \item Check braces.（检查花括号是否成对。）
    \item Check begin/end names.（检查环境开始和结束名称。）
    \item Check file paths and capitalization.（检查文件路径与大小写。）
    \item Check required packages.（检查所需宏包。）
    \item Rebuild references.（重新构建交叉引用和文献引用。）
    \item Clean auxiliary files and rebuild.（清理辅助文件后重新构建。）
\end{enumerate}

% errorbox 只排版排错说明；其中的 \texttt 将错误名称显示为等宽文字，不会触发或模拟一次真实错误。
\begin{errorbox}
不要只看最后一行，也不要一次删除或重写大段源码。例如日志中的
\texttt{Undefined control sequence} 只是现象名称；应回到它指向的第一处相关源码，
判断是拼写、定义顺序还是宏包问题。修好第一处后重新编译，后续连锁错误可能会一起消失。
\end{errorbox}

% 两条终端命令放在 codeblock 中，只作为可复制文本显示；反斜杠续行和命令选项都不会在编译教程时执行。
\subsection{清理后重新构建}

交叉引用、目录或文献状态偶尔会因中断构建而过期。先保存所有源文件并正常重新构建；
只有问题仍然存在、且日志提示旧的辅助状态时，才在项目根目录执行完整清理：

\begin{codeblock}
/Library/TeX/texbin/latexmk -C main.tex
/Library/TeX/texbin/latexmk -synctex=1 -interaction=nonstopmode -file-line-error -xelatex main.tex
\end{codeblock}

第一条命令会移除可再生成的辅助文件和旧输出，第二条命令从 \texttt{main.tex} 重新建立
目录、交叉引用、Biber 文献数据、SyncTeX 映射和 PDF。清理不是修复源码错误的替代品：
如果命令拼写、花括号、环境名或路径本身有误，干净构建仍会在同一根因处失败。

% resultbox 用绿色固定标题显示总结；\textbf 只加粗“章末检查”，后面的命令名用 verb 原样显示。
\begin{resultbox}
\textbf{章末检查：}确认项目只有明确的根文件，各章节由 \verb|\input| 载入；
新增命令先定义后使用且名称不冲突；路径从实际构建入口解释；出现问题时先修第一条有效错误，
最后才清理辅助文件并完整重建。
\end{resultbox}
